\section*{Введение}
\addcontentsline{toc}{section}{Введение}

% Обосновывается актуальность разработки решателя задач для выбранной предметной области. Следует показать практическую или исследовательскую проблему, для решения которой требуется интеллектуальная система. Формулируется цель, перечисляются задачи, выполнение которых необходимо для достижения цели:
% - анализ предметной области,
% - проектирование решателя задач,
% - реализация решателя задач.
% Кратко перечисляются используемые методы решения задач: логический вывод, продукционные правила, поиск в пространстве состояний, эвристический поиск, решение задач удовлетворения ограничений, методы машинного обучения или другие методы, применимые к проекту. Указывается, где и каким образом разработанный решатель задач.
В условиях стремительного роста объемов неструктурированной текстовой информации традиционные методы организации баз знаний и корпоративных поисковых систем сталкиваются с серьезными ограничениями. Классические подходы на основе векторного поиска в системах RAG эффективны для локального контекста, однако не способны отражать сложные семантические и логические связи между концепциями из разных источников. В свою очередь, традиционные графы знаний требуют больших трудозатрат на ручное проектирование и актуализацию. 

Актуальность разработки мультиагентного решателя задач для базы знаний обусловлена необходимостью создания интеллектуальных систем, способных автоматически извлекать, нормализовать, структурировать и непременно перестраивать семантический граф знаний по мере поступления новых материалов, обеспечивая прозрачный и объяснимый процесс поиска ответов.

Целью работы является проектирование и реализация пользовательской подсистемы интеллектуального решателя задач на основе графового RAG, обеспечивающей загрузку документов, наглядную навигацию по базе знаний и получение объяснимых ответов.

Для достижения поставленной цели необходимо решить следующие задачи:
\begin{itemize}
    \item провести анализ предметной области, методов извлечения знаний и существующих архитектур графовых систем RAG;
    \item спроектировать пользовательские сценарии и клиентскую архитектуру для работы с иерархической моделью знаний «сообщество -- домен -- концепция -- карточка знания»;
    \item реализовать клиентский веб-интерфейс загрузки документов, просмотра карточек, навигации и визуализации графа, а также диалога с базой знаний;
    \item интегрировать пользовательский интерфейс с REST API, реализовать отображение состояний фоновой обработки, ошибок, источников и путей рассуждения;
    \item протестировать основные пользовательские сценарии и подготовить демонстрацию программного прототипа.
\end{itemize}

Методы решения задач. В работе применяются методы машинного обучения и обработки естественного языка, методы эвристического поиска в пространстве состояний, векторный поиск ближайших соседей, а также методы алгоритмической нормализации и ранжирования узлов по графовым метрикам.

Практическая применимость. Разработанный решатель задач реализован в виде мультиагентной веб-системы, предназначенной для автоматизации аналитической работы с большими массивами документов, построения персональных и корпоративных баз знаний и ускорения поиска с предоставлением точной траектории рассуждения.
